% tip to kill the wrong type
%!TEX program = xelatex
\documentclass[10pt,journal,compsoc]{IEEEtran}
\IEEEoverridecommandlockouts
% The preceding line is only needed to identify funding in the first footnote. If that is unneeded, please comment it out.
\usepackage{cite}
\usepackage{amsmath,amssymb,amsfonts}
\usepackage{algorithmic}
\usepackage{graphicx}
\usepackage{textcomp}
\usepackage{xcolor}
%\usepackage{}
\usepackage{algorithm}
\usepackage{xeCJK}
\def\BibTeX{{\rm B\kern-.05em{\sc i\kern-.025em b}\kern-.08em
    T\kern-.1667em\lower.7ex\hbox{E}\kern-.125emX}}
\usepackage{fancyhdr}

\begin{document}
实验笔记
\section{实验一}
\subsection{思考题}
题目和答案:

1:如何判断动光栅与静光栅的刻痕已平行?

答案：\textit{观察衍射光斑重合情况：将白纸屏置于光电池前，观察通过双光栅后的衍射光斑。如果两个光栅的刻痕平行，则来自动光栅和静光栅的同级衍射光斑应在屏上完全重合。
调节光栅架角度：缓慢转动静光栅架，观察衍射光斑是否随转动而分离或重合。当光斑重合度最佳时，表明两光栅刻痕已平行。
验证拍频波形：轻敲音叉，在示波器上观察拍频信号。若波形清晰、稳定且幅度较大，则说明两光栅平行度良好，拍频信号已形成。}



2:作为外力驱动音叉谐振曲线时，为什么要固定信号功率?

答案:\textit{固定信号功率是为了保持驱动力恒定，确保实验过程中音叉所受外力一致。如果功率变化，驱动力大小也会改变，导致音叉振动幅度发生变化，从而影响谐振曲线的形状和峰值位置。固定功率后，调节频率可观察音叉在不同频率下的响应特性，得到准确的谐振频率和振幅关系，保证实验数据的可比性和准确性。
}

3:本实验测量方法有何优点?测量微振动位移的灵敏度是多少?

答案:\textit{$10^{-7}$量级}

\subsection{使用到的公式}
\begin{align*}
&F_{平均}=\frac{\text{总波形数}}{T/2}\\
&\text{微振幅：}A=\frac{\text{波形数}}{200}\\
\end{align*}


\section{实验二}
\subsection{思考题}
题目和答案：

1: M-Z光纤干涉仪中要得到清晰的干涉条纹，两路光的光强要接近相同，为什么？

答案：\textit{在双光束干涉中，干涉条纹的对比度（可见度）与两束光的振幅（光强）有关。当两束光的光强相等时，干涉条纹的对比度达到最大，此时明暗条纹分明，清晰度最高。如果两束光的光强相差较大，干涉条纹的对比度会下降，导致条纹模糊不清，不利于观察和测量。}

2: 采用的光纤为光纤通讯用单模光纤，但看到的光斑光强分布并不一定是一个单模高斯分布，为什么？

答案：\textit{实验中使用的光纤虽然是单模光纤，但在耦合过程中可能存在高阶模激发，或者光纤存在弯曲、缺陷等导致模式耦合。此外，光源的相干性、耦合条件的不理想以及光纤端面处理不当等因素也可能导致输出光斑偏离理想的高斯分布。实验中观察到的光斑可能是基模与高阶模的叠加，或者是由于多模干涉等现象造成的。}

3: 比较M-Z光纤干涉仪测量变形量与温度的不同。

答案：
\textit{
  温度对于干涉条纹移动的影响远大于形变的。
    }
\section{实验三}
\subsection{思考题}
题目和答案：

1：真空二极管是否符合欧姆定律？解释之？

答案：\textit{不符合。真空二极管的热电子发射电流由理查逊公式决定，与温度T和逸出功有关，不是线性关系。其伏安特性曲线在空间电荷限制区呈现 $I \propto U^{3/2}$ 关系（Child–Langmuir定律），在饱和区电流基本不变，均不满足欧姆定律 $I \propto U$。}

2：理想二极管中两只栅环起什么作用？

答案：\textit{两只栅环作为保护（补偿）电极，与阳极同电位但绝缘。其作用是消除阴极两端的“冷端效应”和电场不均匀的边缘效应，使中间段阴极的电场和温度分布更均匀，从而将发射区域限制在温度均匀的中央部分，近似为无限长圆柱电极模型，提高测量的准确性。}

3：假如二极管有点漏气的话，伏安特性将会发生怎样的变化？也就是说为什么一定要抽真空？为了得到良好的特性，最大容许残余压强可以有多大？

答案：\textit{若漏气，残余气体分子会与发射电子碰撞，导致电子能量损失甚至引起气体电离，使阳极电流异常增大或不稳定，严重时会烧毁阴极。抽真空是为了消除气体分子对电子运动的干扰，保证热电子发射过程纯净。一般要求真空度优于 $10^{-3}$ Pa（约 $10^{-5}$ Torr）。}

4：伏安特性曲线上当阳极电压为零时电流不为零，并且随阴极温度不同而变化，产生这一现象的原因是什么？

答案：\textit{这是因为热电子从阴极逸出时具有一定的初动能，即使阳极电压为零，部分动能较大的电子仍能克服接触电位差到达阳极，形成微小的电流（即零场电流 $I_s$）。该电流随温度升高而增大，因为温度越高，电子平均动能越大，能到达阳极的电子数越多。}

\subsection{使用到的公式}
\begin{align*}
&\text{理查逊发射公式：} I_S = A S T^2 e^{-\frac{e\phi}{kT}} \\
&\text{理查逊直线法：} \lg\left(\frac{I_S}{T^2}\right) = \lg(AS) - \frac{e\phi}{2.303 k} \cdot \frac{1}{T} \\
&\text{肖特基效应修正：} \lg I'_S = \lg I_S + \frac{0.439}{T} \sqrt{U_a} \\
&\text{温度与灯丝电流关系（近似）：} T = 920.0 + 1600 I_f \; (K) \\
&\text{逸出功计算：} W_0 = e\phi = \frac{k \cdot \text{斜率}}{2.303 \times 10^3} \; (eV)
\end{align*}
    
\section{实验四}
\subsection{思考题}
题目和答案：

1：为何实验上无法观察布拉格散射？

答案：\textit{布拉格衍射需要满足严格的条件：超声波频率较高、声光相互作用距离较大，且光束必须以特定的布拉格角入射。实验中若超声波频率不够高、声光作用长度不足，或入射角偏离布拉格角，则无法满足布拉格衍射条件，此时主要观察到喇曼-奈斯型衍射。此外，声光器件的参数（如声光品质因子）和实验装置的调整精度也会影响布拉格衍射的观察。}

2：声光效应有什么应用？

答案：\textit{声光效应在光学和光电子技术中有广泛应用，主要包括：
（1）声光调制器：通过改变超声波功率来调制衍射光强，用于激光调Q、光通信等；
（2）声光偏转器：通过改变超声波频率来改变衍射光方向，用于激光扫描、光存储等；
（3）声光可调谐滤光器：利用声光效应选择特定波长的光，用于光谱分析、激光调谐等；
（4）声光信号处理：如频谱分析、卷积、相关运算等；
（5）光信息传输：利用声光效应实现光载波的调制与解调。}

\subsection{使用到的公式}
\begin{align*}
&\text{喇曼-奈斯衍射角：} \theta_m = \sin^{-1}\left( m \frac{\lambda_0}{\lambda_s} \right) \quad (m = 0, \pm 1, \pm 2, \dots) \\
&\text{布拉格条件：} \sin\theta_B = \frac{\lambda}{2 n \lambda_s} \\
&\text{布拉格偏转角：} \delta = \frac{\lambda}{n \lambda_s} = \frac{\lambda}{n} \cdot \frac{f_s}{v_s} \\
&\text{声速计算：} v_s = \lambda_s f_s \\
&\text{衍射效率（布拉格）：} \eta = \frac{I_1}{I_0} = \sin^2\left( \frac{\Delta \varphi}{2} \right) \\
&\text{其中相位调制量：} \Delta \varphi = \frac{2\pi}{\lambda} \Delta n L
\end{align*}
\section{实验五}
\subsection{思考题}
题目和答案：

1：简述两种测量法拉第旋转角方法的区别是什么。

答案：\textit{消光法是通过旋转检偏器，使透过检偏器的光强最小，从而直接读取旋转角度，操作简单但精度受消光比影响；倍频法利用交变磁场产生倍频信号，通过示波器观察倍频信号等高时的角度，精度高且能消除光源波动影响，更适合动态测量。}

2：简述磁光效应的应用。

答案：\textit{磁光效应在光隔离器、光调制器、电流传感器、磁场测量、光存储（如磁光盘）、激光通讯和激光雷达中均有重要应用。例如法拉第隔离器可防止激光反射损伤光源，磁光调制器可用于光信号调制，磁光传感器可测量强电流或磁场分布。}

\subsection{使用到的公式}
\begin{align*}
&\text{法拉第旋转角：}\theta_{(\lambda)} = V_{(\lambda)} B d \\
&\text{维尔德常数：}V_{(\lambda)} = \frac{\theta_{(\lambda)}}{B d} \\
&\text{马吕斯定律：}I(a) = I_0 \cos^2 a \\
&\text{倍频法输出光强：}I(a+\theta') = \frac{I_0}{2} \left[1 + \cos 2(a+\theta')\right] \\
&\text{消光条件：}a + \theta = 90^\circ
\end{align*}

\section{实验六}
\subsection{思考题}
题目和答案：

1:锥光干涉图产生的原理是什么？

答案:\textit{利用多个角度入射晶体，不同角度的晶体产生不同的光程差。当满足干涉条件时，在正交偏振镜中通过的是等倾叫干涉和等消光在后焦平面上形成了由黑色十字和曲线组成的图案。}
2:简述电光效应的应用。

答案:\textit{光调制，光线偏转。}
\section{实验七}
\subsection{思考题}
题目和答案：

1：为什么半导体激光器的发散角比气体激光器大？

答案：\textit{半导体激光器谐振腔尺寸小（仅几百微米），且波导层厚度与波长相近，导致衍射效应显著；同时平行结面与垂直结面方向结构不对称，造成两个方向发散角不同，垂直结面方向可达30°–40°。}

2：如何判断半导体激光器是否工作在阈值以上？

答案：\textit{观察输出光功率-电流（P-I）曲线，当电流超过阈值时曲线斜率明显增大，同时光谱从宽自发辐射谱变为窄的激射谱线，且偏振度显著提高。}

3：半导体激光器为什么通常输出TE偏振光？

答案：\textit{GaAs等材料解理面对TE模反射率高于TM模，且波导层对TM模吸收更大，因此TE模阈值增益低，更容易起振，并抑制TM模。}

\subsection{使用到的公式}
\begin{align*}
&\text{阈值电流密度：}J_{th} = \frac{8\pi e n^2 \Delta \nu d}{\eta \lambda_0^2} \left( \alpha + \frac{1}{L} \ln \frac{1}{R} \right) \\
&\text{偏振度：}P = \frac{I_{max} - I_{min}}{I_{max} + I_{min}} \\
&\text{发散角近似：}\theta_\parallel \approx \frac{\lambda}{W},\quad \theta_\perp \approx \frac{\lambda}{d} \\
&\text{纵模间隔：}\Delta \lambda = \frac{\lambda^2}{2nL} \\
&\text{光功率-电流关系：}P = \eta_d \frac{h\nu}{e} (I - I_{th})
\end{align*}

\section{实验八}
\subsection{思考题}
题目和答案：

1：什么是高斯光束？它有什么特点？

答案：\textit{高斯光束是一种光强在横截面上呈高斯分布的光束，其特点包括：光束横截面上光强从中心向外按高斯函数衰减；光束在传播过程中束宽按双曲线规律变化；有确定的束腰位置，在束腰处光斑最小；远场发散角恒定；满足ABCD变换定律，通过透镜变换后仍为高斯光束。}

2：如何测量高斯光束的束腰半径？

答案：\textit{可通过CCD相机在不同位置采集光斑图像，测量各位置的光斑半径ω(z)，然后用双曲线拟合法得到束腰半径ω₀和束腰位置z₀。具体公式为：ω²(z) = A + Bz + Cz²，其中ω₀ = √(A - B²/4C)，z₀ = -B/2C。}

3：什么是M²因子？它有何意义？

答案：\textit{M²因子是光束质量评价参数，定义为：实际光束束腰直径×远场发散角与理想高斯光束束腰直径×远场发散角的比值。M²=1为理想高斯光束；M²>1表示光束质量较差。M²因子综合反映了光束的近场和远场特性，是评价激光束相干性和亮度的重要指标。}

4：高斯光束通过透镜变换有哪些规律？

答案：\textit{高斯光束通过薄透镜变换后仍为高斯光束。变换关系由以下公式描述：
\[ z' = f' + \frac{(z-f')f'^2}{(z-f')^2 + (\pi\omega_0^2/\lambda)^2} \]
\[ \omega_0' = \frac{\omega_0}{\sqrt{(1+z/f')^2 + (Z_{01}/f')^2}} \]
其中z为原束腰到透镜距离，$z'$为新束腰到透镜距离，$f'$为透镜焦距，$Z_{01}=\pi\omega_0^2/\lambda$为瑞利长度。}

\subsection{使用到的公式}
\begin{align*}
&\text{高斯光束表达式：}\Psi_{00}(x,y,z)=\frac{c}{\omega(z)}e^{-\frac{r^2}{\omega^2(z)}}e^{-i\left[k\left(z+\frac{r^2}{2R}\right)-\arctan\frac{z}{f}\right]} \\
&\text{束宽公式：}\omega(z)=\omega_0\sqrt{1+\left(\frac{z}{f}\right)^2} \\
&\text{曲率半径：}R(z)=z+\frac{f^2}{z} \\
&\text{瑞利长度：}f=\frac{\pi\omega_0^2}{\lambda} \\
&\text{远场发散角：}\theta=\lim_{z\to\infty}\frac{\omega(z)}{z}=\frac{\lambda}{\pi\omega_0} \\
&\text{M²因子：}M^2=\frac{\pi}{\lambda}\sqrt{AC-B^2/4} \\
&\text{双曲线拟合：}\omega^2(z)=A+Bz+Cz^2 \\
&\text{透镜变换公式：}z'=f'+\frac{(z-f')f'^2}{(z-f')^2+(\pi\omega_0^2/\lambda)^2} \\
&\text{新束腰半径：}\omega_0'=\frac{\omega_0}{\sqrt{(1+z/f')^2+(Z_{01}/f')^2}}
\end{align*}

\section{实验九}
\subsection{思考题}
题目和答案：

1：霍尔电压是怎样形成的？它的极性与磁场和电流方向（或载子浓度）有什么关系？

答案：\textit{霍尔电压是由于运动的载流子在磁场中受洛伦兹力作用发生偏转，在样品两侧积累正负电荷形成横向电场而产生的。霍尔电压的极性取决于载流子类型：n型半导体（电子导电）霍尔电压为负，p型半导体（空穴导电）霍尔电压为正。对于确定的样品，霍尔电压方向与电流和磁场方向满足右手螺旋关系。}

2：如何观察不等位效应？如何消除它？

答案：\textit{在零磁场（B=0）下通电流I，若A、A'电极间有电压，则为不等位效应电压。消除方法：采用对称测量法，通过改变电流和磁场方向测量四组电压值，取代数平均值，可消除不等位效应和大部分热磁效应。}

3：测量过程中哪些量要保持不变？为什么？

答案：\textit{测量$V_H-I_S$曲线时保持磁场B不变；测量$V_H-B$曲线时保持电流$I_S$不变。这是因为霍尔电压$V_H∝I_S×B$，固定一个变量才能研究另一个变量的线性关系。同时应保持环境温度稳定，避免温度变化影响载流子迁移率。}

4：换向开关的作用原理是什么？测量霍尔电压时为什么要接换向开关？

答案：\textit{换向开关用于改变电流和磁场方向，通过对称测量消除副效应。测量时需接换向开关是因为霍尔效应伴随多种副效应（如不等位效应、温差电效应等），仅单次测量得到的是霍尔电压与各种副效应电压的代数和，通过换向对称测量可消除大部分副效应影响。}

5：工作电流是否可以用交流电源（不考虑表头情况）？为什么？

答案：\textit{不可以。因为霍尔效应是载流子在磁场中定向偏转形成稳定电荷分布的过程，交流电方向周期性变化，无法形成稳定的电荷积累和霍尔电压。且交流电的瞬时性会使载流子来不及响应磁场变化，无法产生稳定的霍尔效应。}

\subsection{使用到的公式}
\begin{align*}
&\text{霍尔电压：}V_H = \frac{1}{ne} \cdot \frac{I_S B}{d} = R_H \frac{I_S B}{d} \\
&\text{霍尔系数：}R_H = \frac{V_H d}{I_S B} = \frac{1}{ne} \quad (\text{单位：m}^3/\text{C}) \\
&\text{载流子浓度：}n = \frac{1}{e R_H} \\
&\text{电导率：}\sigma = \frac{I_S l}{V_\sigma S} \quad (S = b \times d) \\
&\text{迁移率：}\mu = \frac{\sigma}{ne} = |R_H| \sigma \\
&\text{对称测量：}V_H = \frac{V_1 - V_2 + V_3 - V_4}{4} \\
&\text{霍尔灵敏度：}K_H = \frac{V_H}{I_S B} \quad (\text{单位：mV/(mA·T)})
\end{align*}
\section{实验十}
\subsection{思考题}
题目和答案：

1：什么是相位匹配？为什么倍频需要相位匹配？

答案：\textit{相位匹配是指基频光与倍频光在晶体中传播时波矢匹配（$\Delta k = k_{2\omega} - 2k_{\omega} = 0$）。若失配，倍频光会因干涉相消而效率极低，匹配后可实现相干叠加，获得高效倍频输出。}

2：KTP晶体为什么适合用于Nd:YVO4激光器的倍频？

答案：\textit{KTP晶体非线性系数大、损伤阈值高、透光范围覆盖1.064μm和0.53μm，且可实现Ⅱ类相位匹配，匹配角适中，易于实现高效倍频转换。}

3：实验中为什么要先用He-Ne激光准直光路？

答案：\textit{808nm LD光不可见，用632.8nm He-Ne激光作为引导光，便于调节光路共轴，确保泵浦光准确聚焦到晶体内部，提高耦合效率。}

\subsection{使用到的公式}
\begin{align*}
&\text{倍频转换效率：}\eta = \frac{I_{2\omega}}{I_{\omega}} = \frac{8\pi^2 d_{eff}^2 L^2}{\varepsilon_0 c n_{\omega}^2 n_{2\omega} \lambda_{\omega}^2} I_{\omega} \sin c^2\left( \frac{\Delta k L}{2} \right) \\
&\text{相位失配：}\Delta k = k_{2\omega} - 2k_{\omega} = \frac{4\pi}{\lambda_{\omega}} (n_{2\omega} - n_{\omega}) \\
&\text{Ⅱ类相位匹配条件：}n_{\omega}^o + n_{\omega}^e = 2n_{2\omega}^e \\
&\text{有效非线性系数：}d_{eff} = d_{24} \sin\theta \sin 2\phi
\end{align*}

\section{实验十一}
\subsection{思考题}
题目和答案：

1：什么是阿贝成像原理？空间滤波的作用是什么？

答案：\textit{阿贝成像原理认为成像过程是两次傅里叶变换：物面光场经透镜到频谱面（第一次变换），频谱面光场再经透镜到像面（第二次变换）。空间滤波是在频谱面放置滤波器，改变频谱成分，从而改变像的结构，可用于去除噪声、提取特征或实现图像编码。}

2：θ调制实验中如何用白光得到彩色图像？

答案：\textit{不同部分用不同取向光栅调制，白光照明下各方向频谱色散成彩色斑，在频谱面用针孔选择特定颜色通过，像面即合成彩色图像。}

3：低通滤波器和高通滤波器分别有什么作用？

答案：\textit{低通滤波器滤除高频成分，保留图像宏观轮廓，可平滑图像、去除噪声；高通滤波器滤除低频成分，增强图像边缘和细节，突出细微结构。}

\subsection{使用到的公式}
\begin{align*}
&\text{二维傅里叶变换：}F(f_x, f_y) = \iint f(x,y) e^{-i2\pi (f_x x + f_y y)} dx dy \\
&\text{频谱面坐标与空间频率关系：}f_x = \frac{x}{\lambda f},\quad f_y = \frac{y}{\lambda f} \\
&\text{光栅衍射角：}\sin\theta_m = m\frac{\lambda}{d} \\
&\text{阿贝成像光强：}I(u,v) = |\mathcal{F}\{ \mathcal{F}\{O(x,y)\} \cdot H(f_x, f_y) \}|^2
\end{align*}
\end{document}

